%=============================================================================
\section{Architecture Overview}
\label{sec:architecture}
%=============================================================================

Zisk is a zero-knowledge virtual machine that proves correct execution
of RISC-V programs.
It decomposes computation into a single \emph{airgroup} named ``Zisk''
containing 21 specialized AIRs (Algebraic Intermediate Representations),
each responsible for a different aspect of execution:
instruction dispatch, memory, arithmetic, bitwise operations,
cryptographic precompiles, and lookup tables.

The AIRs communicate via \emph{buses}---shared lookup and permutation
arguments that enforce consistency between components.
The central interconnect is the \emph{operation bus}: the Main CPU AIR
dispatches operations to coprocessors (Binary, Arith, etc.) via lookup
assumes, and each coprocessor proves correctness via lookup proves on
the same bus.

%-----------------------------------------------------------------------------
\subsection{Airgroup Structure}
\label{sec:airgroup}

All 21 AIRs belong to a single airgroup (``Zisk'').
The aggregation type is \textbf{sum} (logup) at stage~2,
meaning the global constraint checks that the sum of all $\mathrm{gsum}$
boundary values across all AIRs equals zero.

Configuration parameters:
\begin{itemize}[nosep]
  \item Lattice size: $L = 368$ (no elliptic-curve mode).
  \item Transcript arity: $a = 4$ (Poseidon2 sponge width $= 16$).
  \item Number of public inputs: 68
        (4 for \texttt{rom\_root} $+$ 64 for \texttt{inputs}).
  \item Stage-2 challenges: 2 ($\alpha, \gamma$).
  \item Proof values: 2 at stage~1
        (\texttt{enable\_input\_data}, \texttt{enable\_rom\_data}).
\end{itemize}

%-----------------------------------------------------------------------------
\subsection{AIR Inventory}
\label{sec:air-inventory}

\begin{center}
\renewcommand{\arraystretch}{1.3}
\begin{tabular}{@{}rlrll@{}}
  \toprule
  ID & Name & Trace size & Role & Compressor? \\
  \midrule
  0  & Main              & $2^{22}$ & CPU instruction dispatch        & No \\
  1  & Rom               & $2^{22}$ & Program ROM lookup table        & No \\
  2  & Mem               & $2^{22}$ & Main memory (sorted by addr+step) & No \\
  3  & RomData           & $2^{21}$ & Immutable ROM data region       & No \\
  4  & InputData         & $2^{21}$ & Free input data region          & No \\
  5  & MemAlign          & $2^{21}$ & Unaligned memory access logic   & No \\
  6  & MemAlignByte      & $2^{22}$ & Byte-level memory alignment     & No \\
  7  & MemAlignReadByte  & $2^{22}$ & Read-side byte alignment        & No \\
  8  & MemAlignWriteByte & $2^{22}$ & Write-side byte alignment       & No \\
  9  & Arith             & $2^{21}$ & 64-bit multiply/divide          & No \\
  10 & Binary            & $2^{22}$ & Bitwise AND, OR, XOR, comparisons & No \\
  11 & BinaryAdd         & $2^{22}$ & Dedicated 64-bit addition       & No \\
  12 & BinaryExtension   & $2^{22}$ & Shifts, sign-extension          & No \\
  13 & Add256            & $2^{20}$ & 256-bit addition                & No \\
  14 & ArithEq           & $2^{20}$ & 256-bit field arithmetic (secp256k1, BN254) & Yes \\
  15 & ArithEq384        & $2^{20}$ & 384-bit field arithmetic (BLS12-381) & Yes \\
  16 & Keccakf           & $2^{17}$ & Keccak-f[1600] permutation      & Yes \\
  17 & Sha256f           & $2^{18}$ & SHA-256 compression function    & Yes \\
  18 & SpecifiedRanges   & $2^{20}$ & Range check lookup table        & No \\
  19 & VirtualTable0     & $2^{21}$ & Packed lookup tables (7 tables) & No \\
  20 & VirtualTable1     & $2^{21}$ & Packed lookup tables (3 tables) & No \\
  \bottomrule
\end{tabular}
\end{center}

AIRs marked ``Compressor = Yes'' have STARK verifier circuits exceeding
$2^{17}$ rows, requiring the optional Compressor stage in the recursion
pipeline (see the \emph{Recursion Pipeline}).

%-----------------------------------------------------------------------------
\subsection{STARK Parameters per AIR}
\label{sec:air-params}

Each AIR defines a STARK instance with concrete parameters.
All AIRs share: blowup factor $\beta = 2$, Merkle arity $a = 4$,
quotient degree $d = 2$ (except Rom with $d = 1$), and cubic extension
$\Fext$ ($q_{\mathrm{dim}} = 3$).

\begin{center}
\renewcommand{\arraystretch}{1.2}
\begin{tabular}{@{}lrrrrrr@{}}
  \toprule
  Name & $n$ & Witness & Intmd & Const & Eval pts & FRI rounds \\
  \midrule
  Main              & 22 &  38 &  8 &   3 &   61 & 7 \\
  Rom               & 22 &   1 &  2 &   1 &   18 & 7 \\
  Mem               & 22 &  13 &  3 &   2 &   29 & 7 \\
  RomData           & 21 &   6 &  3 &   2 &   19 & 7 \\
  InputData         & 21 &   9 &  6 &   2 &   27 & 7 \\
  MemAlign          & 21 &  29 &  6 &   2 &   59 & 7 \\
  MemAlignByte      & 22 &  16 &  4 &   1 &   25 & 7 \\
  MemAlignReadByte  & 22 &  10 &  3 &   1 &   18 & 7 \\
  MemAlignWriteByte & 22 &  14 &  4 &   1 &   23 & 7 \\
  Arith             & 21 &  44 & 15 &   1 &   64 & 7 \\
  Binary            & 22 &  39 &  5 &   1 &   49 & 7 \\
  BinaryAdd         & 22 &  10 &  3 &   1 &   18 & 7 \\
  BinaryExtension   & 22 &  29 &  6 &   1 &   40 & 7 \\
  Add256            & 20 &  47 & 17 &   1 &   69 & 6 \\
  ArithEq           & 20 &  39 & 14 &   2 &  434 & 6 \\
  ArithEq384        & 20 &  33 & 14 &   2 &  536 & 6 \\
  Keccakf           & 17 & 2137 & 293 & 2 & 4065 & 5 \\
  Sha256f           & 18 & 102 &  7 &   2 & 1265 & 6 \\
  SpecifiedRanges   & 20 &  33 & 17 &  34 &   88 & 6 \\
  VirtualTable0     & 21 &   8 &  5 &  52 &   69 & 7 \\
  VirtualTable1     & 21 &   8 &  5 &  73 &   90 & 7 \\
  \bottomrule
\end{tabular}
\end{center}

\noindent
\textbf{Columns:}
$n$ = trace size exponent ($N = 2^n$),
Witness = stage-1 committed columns,
Intmd = stage-2 intermediate columns,
Const = constant (setup) columns,
Eval pts = evaluation map entries (polynomial openings in the FRI batching),
FRI rounds = number of folding rounds.
The number of FRI queries ranges from 217 (Keccakf) to 232 (ArithEq384),
targeting approximately 100 bits of security.

\paragraph{Notable outliers.}
Keccakf has by far the most columns (2137 witness $+$ 293 intermediate)
because each Keccak-f round requires extensive bit-level decomposition.
ArithEq384 has the most evaluation points (536) due to the large number
of constraint polynomials in its 384-bit modular arithmetic verification.
The lookup table AIRs (SpecifiedRanges, VirtualTable0, VirtualTable1)
have many constant columns (up to 73) because the precomputed table values
are committed as constant polynomials during setup.

%=============================================================================
\section{Bus Architecture}
\label{sec:buses}
%=============================================================================

%-----------------------------------------------------------------------------
\subsection{Bus Types}
\label{sec:bus-types}

Zisk uses three types of inter-AIR communication:

\begin{itemize}[nosep]
  \item \textbf{Lookup (logup sum-check).}
        One side \emph{proves} rows (provides entries with multiplicities),
        the other side \emph{assumes} rows (consumes entries with selectors).
        The logup argument ensures every assumed tuple was proved.

  \item \textbf{Permutation (multiset equality).}
        Two sides prove/assume with $1{:}1$ matching.
        The grand product argument ensures the multisets are identical.

  \item \textbf{Direct update.}
        Degree-0 constraints for segment chaining and global anchors.
        Used for continuation across execution segments.
\end{itemize}

%-----------------------------------------------------------------------------
\subsection{Bus Inventory}
\label{sec:bus-inventory}

\begin{center}
\renewcommand{\arraystretch}{1.3}
\begin{tabular}{@{}rlll@{}}
  \toprule
  ID & Name & Type & Description \\
  \midrule
  5000 & Operation Bus & Lookup & CPU $\leftrightarrow$ coprocessors \\
  7890 & ROM Bus & Lookup & Main $\leftrightarrow$ Rom \\
  10   & Memory Bus & Permutation & Memory subsystem interconnect \\
  88   & Dual Byte Table & Lookup & Byte-pair decomposition \\
  125  & Binary Table & Lookup & Binary operation lookup \\
  124  & Binary Extension Table & Lookup & Shift/extension lookup \\
  126  & Keccakf Table & Lookup & Keccak round constant lookup \\
  133  & MemAlign ROM & Lookup & Memory alignment microcode \\
  330  & Arith Range Table & Lookup & Arithmetic range check \\
  331  & Arith Table & Lookup & Arithmetic operation lookup \\
  5002 & ArithEq Lt Table & Lookup & Field comparison lookup \\
  5010 & Arith Frops Table & Lookup & Arith frequent operations \\
  5011 & Binary Frops Table & Lookup & Binary frequent operations \\
  5012 & BinaryExt Frops Table & Lookup & BinaryExt frequent operations \\
  \bottomrule
\end{tabular}
\end{center}

%-----------------------------------------------------------------------------
\subsection{Bus Interconnection Summary}
\label{sec:bus-diagram}

The central bus topology is:

\begin{enumerate}[nosep]
  \item \textbf{Main CPU} assumes operations on the \emph{Operation Bus} (5000).
        Every coprocessor (Binary, BinaryAdd, BinaryExtension, Arith,
        Add256, ArithEq, ArithEq384, Keccakf, Sha256f) proves on this bus.

  \item \textbf{Main CPU} assumes instruction fetches on the \emph{ROM Bus} (7890).
        The \textbf{Rom} AIR proves instruction tuples.
        \textbf{RomData} also assumes on the ROM Bus for data reads.

  \item \textbf{Memory operations} use the \emph{Memory Bus} (10) with permutation arguments.
        Main, Mem, RomData, InputData, MemAlign, and the MemAlign byte variants
        all participate.

  \item \textbf{Lookup tables} (SpecifiedRanges, VirtualTable0, VirtualTable1)
        prove on their respective table buses;
        computation AIRs assume range checks and operation tables.
\end{enumerate}

%=============================================================================
\section{CPU: Main AIR}
\label{sec:main}
%=============================================================================

The Main AIR ($2^{22}$ rows) is the central instruction dispatch unit.

%-----------------------------------------------------------------------------
\subsection{Instruction Execution}
\label{sec:main-execution}

Each row of the Main trace represents one clock cycle.
The Main AIR decodes the current instruction and dispatches the operation
to the appropriate coprocessor via the Operation Bus.
It does not perform computation directly---all arithmetic, bitwise,
and cryptographic operations are delegated.

%-----------------------------------------------------------------------------
\subsection{Bus Interactions}
\label{sec:main-buses}

\begin{itemize}[nosep]
  \item \textbf{Operation Bus (5000), assumes.}
        Dispatches operations with the tuple
        $(\mathrm{opcode},\; \mathrm{step},\; 0,\; \mathrm{addr},\; 0,\; \ldots)$.
        The coprocessor that handles this opcode proves the matching tuple.

  \item \textbf{ROM Bus (7890), assumes.}
        Fetches instruction words from the Rom AIR.
        The tuple includes program counter, immediate values,
        and operand encodings.

  \item \textbf{Memory Bus (10), permutation.}
        Issues memory loads and stores.
        The Mem AIR proves the corresponding memory operations.

  \item \textbf{Operation Bus (5000), direct update.}
        Provides continuation anchoring: the initial state (boot address)
        and final state (end address) are published as direct global updates,
        checked by the global constraints.
\end{itemize}

%=============================================================================
\section{Memory Subsystem}
\label{sec:memory}
%=============================================================================

%-----------------------------------------------------------------------------
\subsection{Mem AIR}
\label{sec:mem}

The Mem AIR ($2^{22}$ rows) implements the main read/write memory.
Memory accesses are sorted by $(address, step)$ pairs,
enabling transition constraints that verify temporal consistency:
\begin{itemize}[nosep]
  \item A write at step $t$ stores a value.
  \item A read at step $t' > t$ at the same address returns the most recently written value.
\end{itemize}

Mem supports both 1-byte and 8-byte (aligned) accesses.
The memory bus uses a permutation argument (bus ID~10) to match
loads/stores from the Main AIR and MemAlign subsystem.

%-----------------------------------------------------------------------------
\subsection{RomData AIR}
\label{sec:romdata}

The RomData AIR ($2^{21}$ rows) holds immutable program data.
It follows a \emph{first-write-then-read-only} pattern:
\begin{itemize}[nosep]
  \item Data is written once during initialization.
  \item All subsequent accesses are reads that return the initialized value.
  \item Proved via both the Memory Bus (10) and ROM Bus (7890).
\end{itemize}

Gated by the \texttt{enable\_rom\_data} proof value.

%-----------------------------------------------------------------------------
\subsection{InputData AIR}
\label{sec:inputdata}

The InputData AIR ($2^{21}$ rows) provides free input to the computation.
It follows a \emph{first-read-determines-value} pattern:
\begin{itemize}[nosep]
  \item The first read at an address determines the value (prover's free input).
  \item Subsequent reads return the same value.
  \item Connected via the Memory Bus (10).
\end{itemize}

Gated by the \texttt{enable\_input\_data} proof value.

%-----------------------------------------------------------------------------
\subsection{MemAlign Subsystem}
\label{sec:memalign}

The MemAlign AIR ($2^{21}$ rows) and its three byte-level variants
(MemAlignByte, MemAlignReadByte, MemAlignWriteByte, each $2^{22}$ rows)
form a microprocessor for unaligned memory accesses.

When the Main AIR encounters a memory access that is not naturally aligned,
it delegates to the MemAlign subsystem, which:
\begin{enumerate}[nosep]
  \item Decomposes the unaligned access into aligned sub-accesses.
  \item Performs byte-level read and write operations.
  \item Reassembles the result for the Main AIR.
\end{enumerate}

The MemAlign ROM (bus ID~133) provides microcode for the decomposition.

%=============================================================================
\section{Computation Coprocessors}
\label{sec:coprocessors}
%=============================================================================

Each coprocessor proves operations on the Operation Bus (5000).
The Main AIR assumes tuples of the form
$(\mathrm{opcode}, a_1, a_2, \ldots, a_k)$,
and the coprocessor proves matching tuples with the computed result.

%-----------------------------------------------------------------------------
\subsection{Binary}
\label{sec:binary}

The Binary AIR ($2^{22}$ rows) handles bitwise operations:
AND, OR, XOR, and unsigned/signed comparisons (LTU, LT).
Operations are decomposed into byte-level lookups via the
Binary Table (bus ID~125).

%-----------------------------------------------------------------------------
\subsection{BinaryAdd}
\label{sec:binaryadd}

The BinaryAdd AIR ($2^{22}$ rows) provides dedicated 64-bit addition.
Separated from the Binary AIR for constraint degree optimization.

%-----------------------------------------------------------------------------
\subsection{BinaryExtension}
\label{sec:binaryext}

The BinaryExtension AIR ($2^{22}$ rows) handles shifts
(logical left/right, arithmetic right) and sign-extension operations.
Uses the Binary Extension Table (bus ID~124) for byte-level decomposition.

%-----------------------------------------------------------------------------
\subsection{Arith}
\label{sec:arith}

The Arith AIR ($2^{21}$ rows) handles 64-bit multiplication and division.
Operands are decomposed into 16-bit chunks, and the multiplication
is verified via a schoolbook decomposition with carry propagation.
Uses the Arith Table (bus ID~331) and Arith Range Table (bus ID~330)
for chunk-level lookups.

%=============================================================================
\section{Precompiles}
\label{sec:precompiles}
%=============================================================================

Precompile AIRs provide accelerated circuits for cryptographic operations.
All four precompile AIRs require the Compressor stage in the recursion
pipeline due to their large verifier circuits.

%-----------------------------------------------------------------------------
\subsection{Add256}
\label{sec:add256}

The Add256 AIR ($2^{20}$ rows) provides 256-bit unsigned addition.
Used as a building block for multi-precision arithmetic.

%-----------------------------------------------------------------------------
\subsection{ArithEq}
\label{sec:aritheq}

The ArithEq AIR ($2^{20}$ rows) provides 256-bit modular arithmetic
for elliptic curve operations over secp256k1 and BN254.
Supports:
\begin{itemize}[nosep]
  \item 256-bit modular addition, subtraction, and multiplication.
  \item secp256k1 point addition and doubling.
  \item BN254 curve operations (addition, doubling, complex field arithmetic).
\end{itemize}

Uses the ArithEq Lt Table (bus ID~5002) for field comparison lookups.

%-----------------------------------------------------------------------------
\subsection{ArithEq384}
\label{sec:aritheq384}

The ArithEq384 AIR ($2^{20}$ rows) extends ArithEq to 384-bit fields,
supporting BLS12-381 elliptic curve operations.

%-----------------------------------------------------------------------------
\subsection{Keccakf}
\label{sec:keccakf}

The Keccakf AIR ($2^{17}$ rows) implements the Keccak-f[1600] permutation,
the core of the Keccak/SHA-3 hash function.
Each invocation processes a 1600-bit state through 24 rounds.
Memory operations for reading/writing the state use the Memory Bus (10).
Uses the Keccakf Table (bus ID~126) for round constant lookups.

%-----------------------------------------------------------------------------
\subsection{Sha256f}
\label{sec:sha256f}

The Sha256f AIR ($2^{18}$ rows) implements the SHA-256 compression function.
Each invocation processes a 512-bit message block with a 256-bit state.
Memory operations use the Memory Bus (10).

%=============================================================================
\section{Lookup Tables}
\label{sec:tables}
%=============================================================================

%-----------------------------------------------------------------------------
\subsection{SpecifiedRanges}
\label{sec:specified-ranges}

The SpecifiedRanges AIR ($2^{20}$ rows) provides range check lookups.
It proves that values fall within specified ranges,
used by various AIRs for constraint verification.

%-----------------------------------------------------------------------------
\subsection{VirtualTable0}
\label{sec:vt0}

The VirtualTable0 AIR ($2^{21}$ rows) packs 7 lookup tables into a
single AIR for efficiency:

\begin{center}
\renewcommand{\arraystretch}{1.2}
\begin{tabular}{@{}rll@{}}
  \toprule
  Bus ID & Table & Consumer \\
  \midrule
  331  & Arith Table          & Arith \\
  330  & Arith Range Table    & Arith \\
  5002 & ArithEq Lt Table     & ArithEq, ArithEq384 \\
  124  & Binary Extension Table & BinaryExtension \\
  125  & Binary Table         & Binary \\
  133  & MemAlign ROM         & MemAlign \\
  126  & Keccakf Table        & Keccakf \\
  \bottomrule
\end{tabular}
\end{center}

%-----------------------------------------------------------------------------
\subsection{VirtualTable1}
\label{sec:vt1}

The VirtualTable1 AIR ($2^{21}$ rows) packs 3 frequent-operation
tables:

\begin{center}
\renewcommand{\arraystretch}{1.2}
\begin{tabular}{@{}rll@{}}
  \toprule
  Bus ID & Table & Consumer \\
  \midrule
  5010 & Arith Frops Table        & Arith \\
  5011 & Binary Frops Table       & Binary \\
  5012 & BinaryExt Frops Table    & BinaryExtension \\
  \bottomrule
\end{tabular}
\end{center}

These ``frequent operations'' tables precompute common sub-expressions
used by the computation coprocessors to reduce constraint degree.

%=============================================================================
\section{Global Constraints}
\label{sec:zisk-global-constraints}
%=============================================================================

After all per-AIR proofs are generated, the system must verify
\emph{cross-AIR consistency}: that buses balance, that the execution
begins and ends at the correct state, and that public inputs match.
This verification is performed by the VadcopFinal circuit
(see the \emph{Recursion Pipeline}), which encodes the
global constraints as a single compiled constraint expression.

%-----------------------------------------------------------------------------
\subsection{Bus Balance Equations}
\label{sec:bus-balance}

Each bus connects two or more AIRs via a logup or permutation argument.
Each participating AIR accumulates a running sum ($\mathrm{gsum}$) or
running product ($\mathrm{gprod}$) during its own STARK proof.
The \emph{boundary value} $\mathrm{gsum}^{(a)}[N^{(a)}-1]$
(the final row of the running sum for AIR~$a$) is exported as an
\emph{airgroup value} and carried through the Recursive2 aggregation tree.

For each bus, the global constraint checks that the boundary values
from all participating AIRs cancel:

\begin{itemize}[nosep]
  \item \textbf{Lookup buses} (Operation Bus, ROM Bus, all table buses):
        \[
          \sum_{a \in \mathrm{bus}(b)}
          \mathrm{gsum}^{(a)}[N^{(a)} - 1] = 0.
        \]
        Providers contribute positive multiplicities and consumers
        contribute negative selectors; if every consumed tuple was
        provided, the sum is zero.

  \item \textbf{Permutation buses} (Memory Bus):
        \[
          \prod_{a \in \mathrm{bus}(b)}
          \mathrm{gprod}^{(a)}[N^{(a)} - 1] = 1.
        \]
        The two sides form identical multisets;
        the grand product ratio equals one.
\end{itemize}

\noindent
Concretely, Zisk has 14 buses (\Cref{sec:bus-inventory}),
producing 14 balance equations that the global constraint must enforce.
The single compiled global constraint expression (53{,}890 bytecode
operations) evaluates all 14 balance equations simultaneously,
verifying the entire bus fabric in one check.

%-----------------------------------------------------------------------------
\subsection{Continuation Anchoring}
\label{sec:continuation}

Direct-update constraints anchor the execution to known boundary values:

\begin{itemize}[nosep]
  \item \textbf{Main continuation.}
        The initial program counter equals the boot address.
        The final program counter equals the designated end address.
        Published via direct global updates on the Operation Bus.

  \item \textbf{Memory continuation.}
        Memory base addresses and final states are anchored.

  \item \textbf{Public output.}
        The 64 public input/output values (\texttt{inputs[0..63]}) are
        published via direct global updates, binding the execution trace
        to the claimed public inputs.
\end{itemize}

%-----------------------------------------------------------------------------
\subsection{Airgroup Value Aggregation}
\label{sec:agg-values}

The single airgroup uses sum-type aggregation at stage~2.
During recursive aggregation (see \emph{Recursion Pipeline}, Recursive2),
per-AIR $\mathrm{gsum}$ boundary values are added together.
VadcopFinal checks that the final aggregated sum equals zero
for each bus, confirming cross-AIR consistency.

The aggregation flow is:
\begin{enumerate}[nosep]
  \item Each AIR computes $\mathrm{gsum}^{(a)}[N-1] \in \Fext$
        as a boundary value of its running logup sum.
  \item Recursive2 adds boundary values from child proofs:
        \[
          \mathrm{gsum}_{\mathrm{agg}} = \sum_{a \in \mathrm{children}}
          \mathrm{gsum}^{(a)}.
        \]
  \item VadcopFinal receives the fully aggregated sum and verifies:
        \[
          \mathrm{gsum}_{\mathrm{agg}} = 0 \quad\text{(for each bus)}.
        \]
\end{enumerate}
